% Compact book-start topic list. Keep this student-facing: no source file paths.
\chapter*{অধ্যায়ভিত্তিক বিষয়তালিকা}
\addcontentsline{toc}{chapter}{অধ্যায়ভিত্তিক বিষয়তালিকা}

\begin{compacttable}
\begin{tabular}{|p{0.10\textwidth}|p{0.32\textwidth}|p{0.48\textwidth}|}
\hline
\textbf{অধ্যায়} & \textbf{শিরোনাম} & \textbf{প্রধান আলোচ্য বিষয়} \\
\hline
১ & বিশ্ব ও বাংলাদেশ প্রেক্ষিত & বিশ্বগ্রাম; বিশ্বগ্রামের উপাদান; ভার্চুয়াল রিয়েলিটি; কৃত্রিম বুদ্ধিমত্তা; রোবটিক্স; ক্রায়োসার্জারি; বায়োমেট্রিক্স; বায়োইনফরমেটিক্স; জেনেটিক ইঞ্জিনিয়ারিং; ন্যানোটেকনোলজি; নৈতিকতা; সমাজ ও অর্থনীতিতে আইসিটি \\
\hline
২ & কমিউনিকেশন সিস্টেমস ও নেটওয়ার্কিং & যোগাযোগ ব্যবস্থা; সিগন্যাল ও ব্যান্ডউইথ; ডেটা ট্রান্সমিশন পদ্ধতি ও মোড; গাইডেড ও আনগাইডেড মিডিয়া; মোবাইল যোগাযোগ; নেটওয়ার্কের ধরন; নেটওয়ার্ক ডিভাইস; টপোলজি; ক্লাউড কম্পিউটিং \\
\hline
৩ & সংখ্যা পদ্ধতি ও ডিজিটাল ডিভাইস & সংখ্যা পদ্ধতি; রূপান্তর; বাইনারি যোগ-বিয়োগ; কমপ্লিমেন্ট; কোড; বুলিয়ান অ্যালজেবরা; লজিক গেট; এনকোডার-ডিকোডার; অ্যাডার; রেজিস্টার; কাউন্টার \\
\hline
৪ & ওয়েব ডিজাইন পরিচিতি ও HTML & ওয়েবসাইটের ধারণা; WWW, ব্রাউজার ও সার্চ; URL, HTTP ও ডোমেইন; HTML কাঠামো; ট্যাগ সিনট্যাক্স; ফরম্যাটিং; তালিকা; হাইপারলিংক; ছবি; টেবিল; ফর্ম; div/span/style; ওয়েবসাইট প্রকাশ \\
\hline
৫ & প্রোগ্রামিং ভাষা & প্রোগ্রামের ধারণা; ভাষার প্রজন্ম; অনুবাদক; ডিজাইন ও ডেভেলপমেন্ট; কম্পিউটেশনাল থিংকিং; IPO; ড্রাই রান; অ্যালগরিদম; ফ্লোচার্ট; Python; ডেটা টাইপ; ভ্যারিয়েবল; ইনপুট-আউটপুট; অপারেটর; শর্ত; লুপ; ডেটা স্ট্রাকচার; ফাংশন; কন্ট্রাক্ট; স্কোপ; লোকাল-গ্লোবাল; ফাইল \\
\hline
৬ & ডেটাবেজ ম্যানেজমেন্ট সিস্টেম & ডেটাবেজের মৌলিক ধারণা; DBMS ও RDBMS; টেবিল, ফিল্ড ও রেকর্ড; ডেটা টাইপ; কুয়েরি; সর্টিং; ইনডেক্সিং; রিলেশন; ফর্ম; রিপোর্ট; কর্পোরেট ও সরকারি ডেটাবেজ; নিরাপত্তা; এনক্রিপশন \\
\hline
\end{tabular}
\end{compacttable}
